\subsection{Páginas de progreso y marcas personales}
Se implementaron dos páginas en HTML, CSS y JavaScript, servidas por Express desde el mismo origen que la API. Se eligieron módulos compartidos para la sesión, las solicitudes y los estilos; así se conserva una navegación consistente sin incorporar un framework adicional. Las páginas consultan el perfil autenticado y los endpoints personales; no incluyen resultados simulados en el código de la interfaz.

La página de progreso consulta la racha y los totales por categoría del servidor, y presenta las asistencias confirmadas en un calendario semanal con navegación. Un gráfico circular muestra la distribución por categoría y una franja resume las dos marcas más recientes. También muestra las asistencias recientes y las próximas reservas. La página de marcas permite crear, consultar, modificar y eliminar registros, además de buscar, filtrar y ordenar el historial. Su gráfico reúne únicamente registros del mismo ejercicio, tipo y unidad para evitar comparaciones incompatibles.

Cada operación muestra carga, error o confirmación según la respuesta de la API. Las listas sin registros tienen un estado vacío. Los formularios conservan sus datos ante un rechazo y las eliminaciones solicitan confirmación. Al terminar la sesión se retiran los datos de la pantalla. La comprobación de permisos se mantiene en el servidor.

La interfaz se adaptó a las capturas de Figma aportadas por la usuaria, conservando el encabezado amarillo, la navegación negra y los paneles claros. El acceso utiliza una tarjeta blanca con borde crema, una fotografía alternativa de gimnasio, confirmación de contraseña y control de visibilidad. Se comparó la interfaz con las capturas; no se pudo inspeccionar el archivo editable, por lo que no se declara una reproducción exacta de sus medidas. La revisión se realizó sobre una cuenta y una base sintéticas en navegador de escritorio. La prueba en una tablet física y la integración con las demás pantallas siguen pendientes.
